% Generated human-review packet template.  Source records, Lean statements,
% and raw-source-to-Spec screening fields are substituted by
% scripts/review_dashboard_packet.py.
% Do not hand-edit generated packets; annotate the PDF or reconcile notes into
% the paper-local human-review log.
\documentclass[10pt]{article}
\usepackage[margin=0.43in]{geometry}
\usepackage{fontspec}
\setmainfont{Latin Modern Roman}
\setmonofont{DejaVu Sans Mono}
\usepackage{hyperref}
\usepackage{amssymb}
\usepackage{fvextra}
\usepackage{enumitem}
\setlength{\parindent}{0pt}
\setlength{\parskip}{0.15em}
\linespread{0.92}
\hypersetup{colorlinks=true,linkcolor=blue,urlcolor=blue}
\DefineVerbatimEnvironment{ReviewVerbatim}{Verbatim}{breaklines=true,breakanywhere=true,fontsize=\scriptsize}
\newcommand{\reviewerbox}{%
  \par\noindent\fbox{\parbox{0.96\linewidth}{\rule{0pt}{0.38in}}}\par
}
\newcommand{\reviewmatch}[1]{%
  \par\noindent\textbf{Reviewer check:} $\square$\; #1\par
  \noindent\emph{If left unchecked, explain the discrepancy or uncertainty below.}\par
}

\begin{document}
\fontsize{9}{10}\selectfont

\begin{center}
{\LARGE Human review packet: PKG25NoFreeLunch}\\[0.35em]
\small Generated from byte-pinned source excerpts, the expanded PaperInterface
specifications, and hash-bound raw-source-to-Spec screening records.
\end{center}

\noindent\textbf{Paper:} PKG25NoFreeLunch\\
\textbf{Paper id:} \texttt{PKG25NoFreeLunch}\\
\textbf{Source version:} AAAI 2025, cross-checked against arXiv 2411.15230 source archive SHA256 da7868015db272f3c61cdc88c4e8492b8888f8363332215aa59c509471831c56\\
\textbf{Source record:} \texttt{audit/paper\_statement\_map.json}\\
\textbf{Rows:} 20\\
\textbf{Generated:} 2026-08-18\\
\textbf{Source URL:} \href{https://ojs.aaai.org/index.php/AAAI/article/view/33574/35729}{official source}

\noindent\fbox{\parbox{0.96\linewidth}{\textbf{Review status:} 0/20 source-claim screens. This is a review worksheet while those direct checks remain pending; it is not a final validation receipt.}}\par



\section*{How to use this packet}
For each source claim, compare the exact source excerpts with the expanded
PaperInterface specification. Mark up this PDF or fill the blank reviewer box in the TeX source;
return the annotated file for reconciliation into the paper-local human-review
log.

\section*{Open the interactive dashboard}
The interactive dashboard is an optional alternative to using this PDF; it
presents the same review surface rather than an additional review requirement.
After forking and cloning (or pulling) EconCSLib, run the following from the
repository root. The cache command is needed only the first time in a new clone.
\begin{ReviewVerbatim}
cd <your-EconCSLib-checkout>
lake exe cache get
python3 scripts/review_dashboard.py --paper PKG25NoFreeLunch --serve
\end{ReviewVerbatim}
Open \url{http://127.0.0.1:8765/} in a browser and keep that terminal running
while reviewing. The dashboard records saved annotations in this paper's
reviewer-owned review record.

\section*{Contents}
\begin{itemize}[leftmargin=1.2em]
\item \hyperlink{source-section-f10b5f68e4acf3dc}{Source claims (20)}
\begin{itemize}[leftmargin=1.2em]
\item \hyperlink{source-claim-9384ebab7398c5e4}{source\_definition\_predictorSpec}
\item \hyperlink{source-claim-05c0eaac9818a586}{source\_model\_event\_calibration\_conventionSpec}
\item \hyperlink{source-claim-6ba03b9e70eacd50}{source\_definition\_collaboration\_settingSpec}
\item \hyperlink{source-claim-b8bc36d9b34f3b29}{source\_definition\_rounding\_classifierSpec}
\item \hyperlink{source-claim-9ca4c6ad86e89895}{source\_formula\_individual\_accuracySpec}
\item \hyperlink{source-claim-ca1c9f8cb9fbae0b}{source\_definition\_collaboration\_strategySpec}
\item \hyperlink{source-claim-075776b0c6d0a991}{source\_formula\_induced\_collaboration\_classifierSpec}
\item \hyperlink{source-claim-825d56b262a72a57}{source\_definition\_strategy\_accuracySpec}
\item \hyperlink{source-claim-8acfa10d61ca8cac}{source\_definition\_non\_collaborationSpec}
\item \hyperlink{source-claim-bf0315dd1a18c70e}{source\_theorem1\_no\_free\_lunchSpec}
\item \hyperlink{source-claim-76d57f8e8060f2fd}{source\_definition\_correctSpec}
\item \hyperlink{source-claim-c2574ffc8255a71d}{source\_definition\_incorrectSpec}
\item \hyperlink{source-claim-df70754e8844c37d}{source\_definition\_agreeSpec}
\item \hyperlink{source-claim-b15831d5b4ce3fd5}{source\_definition\_disagreeSpec}
\item \hyperlink{source-claim-bee3bb1dc601f6f9}{source\_proposition6\_linear\_combinationSpec}
\item \hyperlink{source-claim-ca52c024caeea892}{source\_partition\_finite\_measurable\_repairSpec}
\item \hyperlink{source-claim-1b69269ab69dfcc0}{source\_unnumbered\_iff\_restatementSpec}
\item \hyperlink{source-claim-f706fa3b58e00679}{source\_formula\_proposition7\_average\_mixtureSpec}
\item \hyperlink{source-claim-2c1cc7b2a1433aaa}{source\_formula\_proposition9\_s1\_accuraciesSpec}
\item \hyperlink{source-claim-6680595ed06f984e}{source\_formula\_proposition9\_final\_mixtureSpec}
\end{itemize}
\end{itemize}



\clearpage
\hypertarget{source-claim-9384ebab7398c5e4}{}
\hypertarget{source-section-f10b5f68e4acf3dc}{}
\section*{Source claims}
\subsection*{Review row 1}
\textbf{Source locator:} {\footnotesize\raggedright cited publication, lines 7-\allowbreak{}11\par}
\paragraph{Verbatim source input}
\begin{ReviewVerbatim}
A \vocab{predictor} is a function $P: \XX\rightarrow [0,1]$. A predictor $P$ is \vocab{calibrated} on $\DD$ if
\begin{equation}
    \Pr_{(X,Y)\sim \DD}[Y=1\,|\,P(X)=p] = p
\end{equation}
for all $p\in [0,1]$ (more precisely, for all $p\in \text{Image}(P)$).
\end{ReviewVerbatim}


\paragraph{Expanded PaperInterface specification}
\begin{ReviewVerbatim}
∀ {n : ℕ} (S : PKG25NoFreeLunch.JointLawCollaborationSetting n) (i : Fin n) (x : S.X), 0 ≤ S.pred i x ∧ S.pred i x ≤ 1
\end{ReviewVerbatim}

\paragraph{Lean proof endpoint (built separately)}\begin{ReviewVerbatim}
PKG25NoFreeLunch.source_definition_predictorSpec_proof
\end{ReviewVerbatim}

\paragraph{Recorded source-to-Spec screening}
\textbf{Verdict:} \texttt{not recorded}
\\
\textbf{Reason:} not recorded
\reviewmatch{Matches source input}
\noindent\textbf{Reviewer annotation}\par
\reviewerbox
\clearpage
\hypertarget{source-claim-05c0eaac9818a586}{}
\section*{Review row 2}
\textbf{Source locator:} {\footnotesize\raggedright cited publication, lines 7-\allowbreak{}11\par}
\paragraph{Verbatim source input}
\begin{ReviewVerbatim}
A \vocab{predictor} is a function $P: \XX\rightarrow [0,1]$. A predictor $P$ is \vocab{calibrated} on $\DD$ if
\begin{equation}
    \Pr_{(X,Y)\sim \DD}[Y=1\,|\,P(X)=p] = p
\end{equation}
for all $p\in [0,1]$ (more precisely, for all $p\in \text{Image}(P)$).
\end{ReviewVerbatim}


\paragraph{Expanded PaperInterface specification}
\begin{ReviewVerbatim}
∀ {n : ℕ} (S : PKG25NoFreeLunch.JointLawCollaborationSetting n) (i : Fin n) (A : Set ℝ),
  MeasurableSet A →
    ∫ (z : S.X × PKG25NoFreeLunch.Label), {z | S.pred i z.1 ∈ A ∧ z.2 = true}.indicator (fun x => 1) z ∂S.joint =
      ∫ (z : S.X × PKG25NoFreeLunch.Label), {z | S.pred i z.1 ∈ A}.indicator (fun z => S.pred i z.1) z ∂S.joint
\end{ReviewVerbatim}

\paragraph{Lean proof endpoint (built separately)}\begin{ReviewVerbatim}
PKG25NoFreeLunch.source_model_event_calibration_conventionSpec_proof
\end{ReviewVerbatim}

\paragraph{Recorded source-to-Spec screening}
\textbf{Verdict:} \texttt{not recorded}
\\
\textbf{Reason:} not recorded
\reviewmatch{Matches source input}
\noindent\textbf{Reviewer annotation}\par
\reviewerbox
\clearpage
\hypertarget{source-claim-6ba03b9e70eacd50}{}
\section*{Review row 3}
\textbf{Source locator:} {\footnotesize\raggedright cited publication, lines 13-\allowbreak{}19\par}
\paragraph{Verbatim source input}
\begin{ReviewVerbatim}
We may now define collaboration settings.
\begin{definition}A \textbf{collaboration setting} is an ordered tuple
\begin{equation}
    S=(\DD, P_1, \cdots, P_n),
\end{equation}
where $\DD$ is a probability distribution over $\XX\times \{0,1\}$ and $P_1,\cdots,P_n$ are each calibrated predictors on $\DD$.
\end{definition}
\end{ReviewVerbatim}


\paragraph{Expanded PaperInterface specification}
\begin{ReviewVerbatim}
∀ (n : ℕ), PKG25NoFreeLunch.source_definition_probability_collaboration_settingSpec n
\end{ReviewVerbatim}

\paragraph{Lean proof endpoint (built separately)}\begin{ReviewVerbatim}
PKG25NoFreeLunch.source_definition_collaboration_settingSpec_proof
\end{ReviewVerbatim}

\paragraph{Recorded source-to-Spec screening}
\textbf{Verdict:} \texttt{not recorded}
\\
\textbf{Reason:} not recorded
\reviewmatch{Matches source input}
\noindent\textbf{Reviewer annotation}\par
\reviewerbox
\clearpage
\hypertarget{source-claim-b8bc36d9b34f3b29}{}
\section*{Review row 4}
\textbf{Source locator:} {\footnotesize\raggedright cited publication, lines 21-\allowbreak{}24\par}
\paragraph{Verbatim source input}
\begin{ReviewVerbatim}
The predictor $P_i$ induces the classifier $\Yhat_i: x\mapsto \round{P_i(x)},$ where $\round{\cdot}$ is the rounding operator (without loss of generality, set $\round{0.5}=1$). $\Yhat_i$ achieves 0-1 accuracy
\begin{equation}
    \acc_i(S) := \EE_{(X,Y)\sim \DD}[\max\{P_i(X), 1 - P_i(X)\}].
\end{equation}
\end{ReviewVerbatim}


\paragraph{Expanded PaperInterface specification}
\begin{ReviewVerbatim}
PKG25NoFreeLunch.source_definition_rounding_conventionSpec ∧
  ∀ {n : ℕ} (S : PKG25NoFreeLunch.JointLawCollaborationSetting n) (i : Fin n) (x : S.X),
    S.agentClassifier i x = PKG25NoFreeLunch.roundProb (S.pred i x)
\end{ReviewVerbatim}

\paragraph{Lean proof endpoint (built separately)}\begin{ReviewVerbatim}
PKG25NoFreeLunch.source_definition_rounding_classifierSpec_proof
\end{ReviewVerbatim}

\paragraph{Recorded source-to-Spec screening}
\textbf{Verdict:} \texttt{not recorded}
\\
\textbf{Reason:} not recorded
\reviewmatch{Matches source input}
\noindent\textbf{Reviewer annotation}\par
\reviewerbox
\clearpage
\hypertarget{source-claim-9ca4c6ad86e89895}{}
\section*{Review row 5}
\textbf{Source locator:} {\footnotesize\raggedright cited publication, lines 21-\allowbreak{}24\par}
\paragraph{Verbatim source input}
\begin{ReviewVerbatim}
The predictor $P_i$ induces the classifier $\Yhat_i: x\mapsto \round{P_i(x)},$ where $\round{\cdot}$ is the rounding operator (without loss of generality, set $\round{0.5}=1$). $\Yhat_i$ achieves 0-1 accuracy
\begin{equation}
    \acc_i(S) := \EE_{(X,Y)\sim \DD}[\max\{P_i(X), 1 - P_i(X)\}].
\end{equation}
\end{ReviewVerbatim}


\paragraph{Expanded PaperInterface specification}
\begin{ReviewVerbatim}
∀ {n : ℕ} (S : PKG25NoFreeLunch.JointLawCollaborationSetting n) (i : Fin n),
  S.agentAccuracy i = ∫ (z : S.X × PKG25NoFreeLunch.Label), max (S.pred i z.1) (1 - S.pred i z.1) ∂S.joint
\end{ReviewVerbatim}

\paragraph{Lean proof endpoint (built separately)}\begin{ReviewVerbatim}
PKG25NoFreeLunch.source_formula_individual_accuracySpec_proof
\end{ReviewVerbatim}

\paragraph{Recorded source-to-Spec screening}
\textbf{Verdict:} \texttt{not recorded}
\\
\textbf{Reason:} not recorded
\reviewmatch{Matches source input}
\noindent\textbf{Reviewer annotation}\par
\reviewerbox
\clearpage
\hypertarget{source-claim-ca1c9f8cb9fbae0b}{}
\section*{Review row 6}
\textbf{Source locator:} {\footnotesize\raggedright cited publication, lines 27-\allowbreak{}35\par}
\paragraph{Verbatim source input}
\begin{ReviewVerbatim}
A collaboration strategy is a way to combine the predicted probabilities of each agent to make a classification.
\begin{definition}
A \textbf{collaboration strategy} is a deterministic function $\CC: [0,1]^n \rightarrow \{0,1\}.$
\end{definition}
Here, $\CC$ should be interpreted as a function that takes in $n$ predicted probabilities and returns a 0-1 classification. Given a collaboration setting $S=(\DD, P_1, \cdots, P_n)$, a collaboration strategy $\CC$ induces the classifier
\begin{equation}
    \Yhat_\CC: \XX \rightarrow \{0,1\},\quad x \mapsto \CC(P_1(x),\cdots,P_n(x)).
\end{equation}
Let $\acc_\CC(S)$ denote the 0-1 accuracy of $\Yhat_\CC$ on $\DD$. In other words, for each $x\in \XX$, each agent $i$ produces a predicted probability $P_i(x).$ The collaboration strategy aggregates these predictions into a binary classification.
\end{ReviewVerbatim}


\paragraph{Expanded PaperInterface specification}
\begin{ReviewVerbatim}
∀ (n : ℕ), PKG25NoFreeLunch.source_definition_collaboration_strategy n = PKG25NoFreeLunch.CollaborationStrategy n
\end{ReviewVerbatim}

\paragraph{Lean proof endpoint (built separately)}\begin{ReviewVerbatim}
PKG25NoFreeLunch.source_definition_collaboration_strategySpec_proof
\end{ReviewVerbatim}

\paragraph{Recorded source-to-Spec screening}
\textbf{Verdict:} \texttt{not recorded}
\\
\textbf{Reason:} not recorded
\reviewmatch{Matches source input}
\noindent\textbf{Reviewer annotation}\par
\reviewerbox
\clearpage
\hypertarget{source-claim-075776b0c6d0a991}{}
\section*{Review row 7}
\textbf{Source locator:} {\footnotesize\raggedright cited publication, lines 31-\allowbreak{}35\par}
\paragraph{Verbatim source input}
\begin{ReviewVerbatim}
Here, $\CC$ should be interpreted as a function that takes in $n$ predicted probabilities and returns a 0-1 classification. Given a collaboration setting $S=(\DD, P_1, \cdots, P_n)$, a collaboration strategy $\CC$ induces the classifier
\begin{equation}
    \Yhat_\CC: \XX \rightarrow \{0,1\},\quad x \mapsto \CC(P_1(x),\cdots,P_n(x)).
\end{equation}
Let $\acc_\CC(S)$ denote the 0-1 accuracy of $\Yhat_\CC$ on $\DD$. In other words, for each $x\in \XX$, each agent $i$ produces a predicted probability $P_i(x).$ The collaboration strategy aggregates these predictions into a binary classification.
\end{ReviewVerbatim}


\paragraph{Expanded PaperInterface specification}
\begin{ReviewVerbatim}
∀ {n : ℕ} (S : PKG25NoFreeLunch.JointLawCollaborationSetting n) (C : PKG25NoFreeLunch.CollaborationStrategy n)
  (x : S.X), S.strategyClassifier C x = C fun i => S.pred i x
\end{ReviewVerbatim}

\paragraph{Lean proof endpoint (built separately)}\begin{ReviewVerbatim}
PKG25NoFreeLunch.source_formula_induced_collaboration_classifierSpec_proof
\end{ReviewVerbatim}

\paragraph{Recorded source-to-Spec screening}
\textbf{Verdict:} \texttt{not recorded}
\\
\textbf{Reason:} not recorded
\reviewmatch{Matches source input}
\noindent\textbf{Reviewer annotation}\par
\reviewerbox
\clearpage
\hypertarget{source-claim-825d56b262a72a57}{}
\section*{Review row 8}
\textbf{Source locator:} {\footnotesize\raggedright cited publication, lines 35\par}
\paragraph{Verbatim source input}
\begin{ReviewVerbatim}
Let $\acc_\CC(S)$ denote the 0-1 accuracy of $\Yhat_\CC$ on $\DD$. In other words, for each $x\in \XX$, each agent $i$ produces a predicted probability $P_i(x).$ The collaboration strategy aggregates these predictions into a binary classification.
\end{ReviewVerbatim}


\paragraph{Expanded PaperInterface specification}
\begin{ReviewVerbatim}
∀ {n : ℕ} (S : PKG25NoFreeLunch.JointLawCollaborationSetting n) (C : PKG25NoFreeLunch.CollaborationStrategy n),
  PKG25NoFreeLunch.source_assumption_strategy_expectation_well_formed S C →
    S.strategyAccuracy C =
      ∫ (z : S.X × PKG25NoFreeLunch.Label), if S.strategyClassifier C z.1 = z.2 then 1 else 0 ∂S.joint
\end{ReviewVerbatim}

\paragraph{Lean proof endpoint (built separately)}\begin{ReviewVerbatim}
PKG25NoFreeLunch.source_definition_strategy_accuracySpec_proof
\end{ReviewVerbatim}

\paragraph{Recorded source-to-Spec screening}
\textbf{Verdict:} \texttt{not recorded}
\\
\textbf{Reason:} not recorded
\reviewmatch{Matches source input}
\noindent\textbf{Reviewer annotation}\par
\reviewerbox
\clearpage
\hypertarget{source-claim-8acfa10d61ca8cac}{}
\section*{Review row 9}
\textbf{Source locator:} {\footnotesize\raggedright cited publication, lines 42-\allowbreak{}54\par}
\paragraph{Verbatim source input}
\begin{ReviewVerbatim}
\begin{definition}\label{def:non-collaborative}
    A collaboration strategy $\CC:[0,1]^n\rightarrow \{0,1\}$ is \textbf{non-collaborative} if there exists $k\in [n]$ and $\alpha\in \{0,1\}$ such that
    \begin{equation}
    \CC(p_1,p_2,\cdots,p_n) =
        \begin{cases}
        \round{p_k} &\quad \text{$p_k\neq \frac{1}{2}$} \\
            \alpha &\quad \text{$p_k=\frac{1}{2}$}
        \end{cases}
    \end{equation}
    for all $(p_1,p_2,\cdots,p_n)\in (0,1)^n$.
\end{definition}

A collaboration strategy is non-collaborative if it essentially always defers to the classification of a single agent $k\in [n]$. The definition admits two exceptions. First, when $p_k = \frac{1}{2}$, the collaboration may select either $0$ or $1$, but must always select the same such value; this choice has no effect on the accuracy of the classifier, so this exception can be essentially ignored. Second, the collaboration strategy need not select $\round{p_k}$ when $(p_1,p_2,\cdots,p_n)\notin (0,1)^n$---i.e., when $p_i\in \{0,1\}$ for some $i\in [n]$; this exception is somewhat more interesting, and we will return to it after stating our main result.
\end{ReviewVerbatim}


\paragraph{Expanded PaperInterface specification}
\begin{ReviewVerbatim}
∀ {n : ℕ} (C : PKG25NoFreeLunch.CollaborationStrategy n),
  PKG25NoFreeLunch.NonCollaborative C ↔
    ∃ k α, PKG25NoFreeLunch.DefersAwayFromHalf C k ∧ PKG25NoFreeLunch.ConstantOnHalfSlice C k α
\end{ReviewVerbatim}

\paragraph{Lean proof endpoint (built separately)}\begin{ReviewVerbatim}
PKG25NoFreeLunch.source_definition_non_collaborationSpec_proof
\end{ReviewVerbatim}

\paragraph{Recorded source-to-Spec screening}
\textbf{Verdict:} \texttt{not recorded}
\\
\textbf{Reason:} not recorded
\reviewmatch{Matches source input}
\noindent\textbf{Reviewer annotation}\par
\reviewerbox
\clearpage
\hypertarget{source-claim-bf0315dd1a18c70e}{}
\section*{Review row 10}
\textbf{Source locator:} {\footnotesize\raggedright cited publication, lines 56-\allowbreak{}63\par}
\paragraph{Verbatim source input}
\begin{ReviewVerbatim}
\setcounter{theorem}{0}
\begin{theorem}\label{thm:main}
    Every reliable collaboration strategy is non-collaborative.




\end{theorem}
\end{ReviewVerbatim}


\paragraph{Expanded PaperInterface specification}
\begin{ReviewVerbatim}
∀ {n : ℕ} [Nonempty (Fin n)] (C : PKG25NoFreeLunch.CollaborationStrategy n),
  PKG25NoFreeLunch.ReliableJointLaw C → PKG25NoFreeLunch.NonCollaborative C
\end{ReviewVerbatim}

\paragraph{Lean proof endpoint (built separately)}\begin{ReviewVerbatim}
PKG25NoFreeLunch.source_theorem1_no_free_lunchSpec_proof
\end{ReviewVerbatim}

\paragraph{Recorded source-to-Spec screening}
\textbf{Verdict:} \texttt{not recorded}
\\
\textbf{Reason:} not recorded
\reviewmatch{Matches source input}
\noindent\textbf{Reviewer annotation}\par
\reviewerbox
\clearpage
\hypertarget{source-claim-76d57f8e8060f2fd}{}
\section*{Review row 11}
\textbf{Source locator:} {\footnotesize\raggedright cited publication, lines 83-\allowbreak{}92\par}
\paragraph{Verbatim source input}
\begin{ReviewVerbatim}
\begin{definition}\label{def:correct}
For a collaboration setting $(\DD, P_1,\cdots,P_n)$ and $x\in \XX,$ we say that
\begin{itemize}
    \item agent $i$ is \textbf{correct} on $x$ if $$\Yhat_i(x) = \round{\Pr_{(X,Y)\sim \DD}[Y=1\,|\,X=x]}$$
    \item agent $i$ is \textbf{incorrect} on $x$ if $$\Yhat_i(x) \neq \round{\Pr_{(X,Y)\sim \DD}[Y=1\,|\,X=x]}$$
    \item agents $i$ and $j$ \textbf{agree} on $x$ if $$\Yhat_i(x) = \Yhat_j(x)$$
    \item agents $i$ and $j$ \textbf{disagree} on $x$ if $$\Yhat_i(x) \neq \Yhat_j(x).$$
\end{itemize}
For each of these statements, we can replace $i$ or $j$ with a collaboration strategy $\CC$.
\end{definition}
\end{ReviewVerbatim}


\paragraph{Expanded PaperInterface specification}
\begin{ReviewVerbatim}
∀ (prediction : PKG25NoFreeLunch.Label) (eta : ℝ),
  PKG25NoFreeLunch.source_definition_correct_on prediction eta ↔ prediction = PKG25NoFreeLunch.roundProb eta
\end{ReviewVerbatim}

\paragraph{Lean proof endpoint (built separately)}\begin{ReviewVerbatim}
PKG25NoFreeLunch.source_definition_correctSpec_proof
\end{ReviewVerbatim}

\paragraph{Recorded source-to-Spec screening}
\textbf{Verdict:} \texttt{not recorded}
\\
\textbf{Reason:} not recorded
\reviewmatch{Matches source input}
\noindent\textbf{Reviewer annotation}\par
\reviewerbox
\clearpage
\hypertarget{source-claim-c2574ffc8255a71d}{}
\section*{Review row 12}
\textbf{Source locator:} {\footnotesize\raggedright cited publication, lines 83-\allowbreak{}92\par}
\paragraph{Verbatim source input}
\begin{ReviewVerbatim}
\begin{definition}\label{def:correct}
For a collaboration setting $(\DD, P_1,\cdots,P_n)$ and $x\in \XX,$ we say that
\begin{itemize}
    \item agent $i$ is \textbf{correct} on $x$ if $$\Yhat_i(x) = \round{\Pr_{(X,Y)\sim \DD}[Y=1\,|\,X=x]}$$
    \item agent $i$ is \textbf{incorrect} on $x$ if $$\Yhat_i(x) \neq \round{\Pr_{(X,Y)\sim \DD}[Y=1\,|\,X=x]}$$
    \item agents $i$ and $j$ \textbf{agree} on $x$ if $$\Yhat_i(x) = \Yhat_j(x)$$
    \item agents $i$ and $j$ \textbf{disagree} on $x$ if $$\Yhat_i(x) \neq \Yhat_j(x).$$
\end{itemize}
For each of these statements, we can replace $i$ or $j$ with a collaboration strategy $\CC$.
\end{definition}
\end{ReviewVerbatim}


\paragraph{Expanded PaperInterface specification}
\begin{ReviewVerbatim}
∀ (prediction : PKG25NoFreeLunch.Label) (eta : ℝ),
  PKG25NoFreeLunch.source_definition_incorrect_on prediction eta ↔ prediction ≠ PKG25NoFreeLunch.roundProb eta
\end{ReviewVerbatim}

\paragraph{Lean proof endpoint (built separately)}\begin{ReviewVerbatim}
PKG25NoFreeLunch.source_definition_incorrectSpec_proof
\end{ReviewVerbatim}

\paragraph{Recorded source-to-Spec screening}
\textbf{Verdict:} \texttt{not recorded}
\\
\textbf{Reason:} not recorded
\reviewmatch{Matches source input}
\noindent\textbf{Reviewer annotation}\par
\reviewerbox
\clearpage
\hypertarget{source-claim-df70754e8844c37d}{}
\section*{Review row 13}
\textbf{Source locator:} {\footnotesize\raggedright cited publication, lines 83-\allowbreak{}92\par}
\paragraph{Verbatim source input}
\begin{ReviewVerbatim}
\begin{definition}\label{def:correct}
For a collaboration setting $(\DD, P_1,\cdots,P_n)$ and $x\in \XX,$ we say that
\begin{itemize}
    \item agent $i$ is \textbf{correct} on $x$ if $$\Yhat_i(x) = \round{\Pr_{(X,Y)\sim \DD}[Y=1\,|\,X=x]}$$
    \item agent $i$ is \textbf{incorrect} on $x$ if $$\Yhat_i(x) \neq \round{\Pr_{(X,Y)\sim \DD}[Y=1\,|\,X=x]}$$
    \item agents $i$ and $j$ \textbf{agree} on $x$ if $$\Yhat_i(x) = \Yhat_j(x)$$
    \item agents $i$ and $j$ \textbf{disagree} on $x$ if $$\Yhat_i(x) \neq \Yhat_j(x).$$
\end{itemize}
For each of these statements, we can replace $i$ or $j$ with a collaboration strategy $\CC$.
\end{definition}
\end{ReviewVerbatim}


\paragraph{Expanded PaperInterface specification}
\begin{ReviewVerbatim}
∀ (prediction1 prediction2 : PKG25NoFreeLunch.Label),
  PKG25NoFreeLunch.source_definition_agree_on prediction1 prediction2 ↔ prediction1 = prediction2
\end{ReviewVerbatim}

\paragraph{Lean proof endpoint (built separately)}\begin{ReviewVerbatim}
PKG25NoFreeLunch.source_definition_agreeSpec_proof
\end{ReviewVerbatim}

\paragraph{Recorded source-to-Spec screening}
\textbf{Verdict:} \texttt{not recorded}
\\
\textbf{Reason:} not recorded
\reviewmatch{Matches source input}
\noindent\textbf{Reviewer annotation}\par
\reviewerbox
\clearpage
\hypertarget{source-claim-b15831d5b4ce3fd5}{}
\section*{Review row 14}
\textbf{Source locator:} {\footnotesize\raggedright cited publication, lines 83-\allowbreak{}92\par}
\paragraph{Verbatim source input}
\begin{ReviewVerbatim}
\begin{definition}\label{def:correct}
For a collaboration setting $(\DD, P_1,\cdots,P_n)$ and $x\in \XX,$ we say that
\begin{itemize}
    \item agent $i$ is \textbf{correct} on $x$ if $$\Yhat_i(x) = \round{\Pr_{(X,Y)\sim \DD}[Y=1\,|\,X=x]}$$
    \item agent $i$ is \textbf{incorrect} on $x$ if $$\Yhat_i(x) \neq \round{\Pr_{(X,Y)\sim \DD}[Y=1\,|\,X=x]}$$
    \item agents $i$ and $j$ \textbf{agree} on $x$ if $$\Yhat_i(x) = \Yhat_j(x)$$
    \item agents $i$ and $j$ \textbf{disagree} on $x$ if $$\Yhat_i(x) \neq \Yhat_j(x).$$
\end{itemize}
For each of these statements, we can replace $i$ or $j$ with a collaboration strategy $\CC$.
\end{definition}
\end{ReviewVerbatim}


\paragraph{Expanded PaperInterface specification}
\begin{ReviewVerbatim}
∀ (prediction1 prediction2 : PKG25NoFreeLunch.Label),
  PKG25NoFreeLunch.source_definition_disagree_on prediction1 prediction2 ↔ prediction1 ≠ prediction2
\end{ReviewVerbatim}

\paragraph{Lean proof endpoint (built separately)}\begin{ReviewVerbatim}
PKG25NoFreeLunch.source_definition_disagreeSpec_proof
\end{ReviewVerbatim}

\paragraph{Recorded source-to-Spec screening}
\textbf{Verdict:} \texttt{not recorded}
\\
\textbf{Reason:} not recorded
\reviewmatch{Matches source input}
\noindent\textbf{Reviewer annotation}\par
\reviewerbox
\clearpage
\hypertarget{source-claim-bee3bb1dc601f6f9}{}
\section*{Review row 15}
\textbf{Source locator:} {\footnotesize\raggedright cited publication, lines 100-\allowbreak{}107\par}
\paragraph{Verbatim source input}
\begin{ReviewVerbatim}
\begin{proposition}\label{prop:linear-combination} \textbf{Linear combinations of settings.}
    Consider $\ell$ collaboration settings $S_1, \cdots, S_\ell$. Then for all $(\lambda_1, \cdots, \lambda_\ell)\in \Delta^\ell$, there exists a collaboration setting $S$ such that
    \begin{align}
        \acc_i(S) &= \sum_{m=1}^\ell \lambda_m \acc_i(S_m)\label{eq:monstera-1}\\
        \acc_\CC(S) &= \sum_{m=1}^\ell \lambda_m \acc_\CC(S_m)\label{eq:monstera-2}
    \end{align}
    for all $i\in [n]$ and collaboration strategies $\CC$.
\end{proposition}
\end{ReviewVerbatim}


\paragraph{Expanded PaperInterface specification}
\begin{ReviewVerbatim}
∀ {n ell : ℕ} (S : Fin ell → PKG25NoFreeLunch.JointLawCollaborationSetting n) (w : Fin ell → ℝ),
  (∀ (r : Fin ell), 0 ≤ w r) →
    ∑ r, w r = 1 →
      ∃ Smix,
        (∀ (i : Fin n), Smix.agentAccuracy i = ∑ r, w r * (S r).agentAccuracy i) ∧
          ∀ (C : PKG25NoFreeLunch.CollaborationStrategy n),
            (∀ (r : Fin ell), PKG25NoFreeLunch.source_assumption_strategy_expectation_well_formed (S r) C) →
              PKG25NoFreeLunch.source_assumption_strategy_expectation_well_formed Smix C ∧
                Smix.strategyAccuracy C = ∑ r, w r * (S r).strategyAccuracy C
\end{ReviewVerbatim}

\paragraph{Lean proof endpoint (built separately)}\begin{ReviewVerbatim}
PKG25NoFreeLunch.source_proposition6_linear_combinationSpec_proof
\end{ReviewVerbatim}

\paragraph{Recorded source-to-Spec screening}
\textbf{Verdict:} \texttt{not recorded}
\\
\textbf{Reason:} not recorded
\reviewmatch{Matches source input}
\noindent\textbf{Reviewer annotation}\par
\reviewerbox
\clearpage
\hypertarget{source-claim-ca52c024caeea892}{}
\section*{Review row 16}
\textbf{Source locator:} {\footnotesize\raggedright cited publication, lines 127-\allowbreak{}132\par}
\paragraph{Verbatim source input}
\begin{ReviewVerbatim}
\paragraph{Building Calibrated Predictors from Partitions.}
Finally, given a distribution $\DD$ over $\XX\times \{0,1\}$, we show how partitions of the input space $\XX$ induce calibrated predictors. Indeed, let $\AA_i$ be a partition of $\XX$. Then $\AA_i$ induces a calibrated predictor $P_i$, where for $x\in A \in \AA_i$,
\begin{equation}
    P_i(x) := \Pr_{(X,Y)\sim \DD}[Y=1\,|\,X\in A].
\end{equation}
Here $P_i$ is the Bayes-optimal predictor given a ``coarsening'' of the input space into the partitions $\XX.$ In this way, a collaboration setting may also be identified by an ordered tuple $(\DD, \AA_1, \cdots, \AA_n).$ This approach is central to the subsequent proofs.
\end{ReviewVerbatim}


\paragraph{Expanded PaperInterface specification}
\begin{ReviewVerbatim}
(∀ {X : Type u_1} {Cell : Type u_2} [inst : MeasurableSpace X] [inst_1 : MeasurableSpace Cell]
    [MeasurableSingletonClass Cell] (joint : MeasureTheory.Measure (X × PKG25NoFreeLunch.Label))
    [MeasureTheory.IsProbabilityMeasure joint] (cell : X → Cell),
    Measurable cell →
      ∀ (x : X),
        0 < PKG25NoFreeLunch.jointLawPartitionCellMass joint cell (cell x) →
          PKG25NoFreeLunch.jointLawPartitionPredictor joint cell x =
            PKG25NoFreeLunch.jointLawPartitionCellTrueMass joint cell (cell x) /
              PKG25NoFreeLunch.jointLawPartitionCellMass joint cell (cell x)) ∧
  (∀ {X : Type u_3} {Cell : Type u_4} [inst : MeasurableSpace X] [inst_1 : MeasurableSpace Cell]
      [MeasurableSingletonClass Cell] (joint : MeasureTheory.Measure (X × PKG25NoFreeLunch.Label))
      [MeasureTheory.IsProbabilityMeasure joint] (cell : X → Cell),
      Measurable cell →
        ∀ (x : X),
          PKG25NoFreeLunch.jointLawPartitionCellMass joint cell (cell x) = 0 →
            PKG25NoFreeLunch.jointLawPartitionPredictor joint cell x = 0) ∧
    ∀ {X : Type u_5} {Cell : Type u_6} [inst : MeasurableSpace X] [Fintype Cell] [inst_2 : MeasurableSpace Cell]
      [MeasurableSingletonClass Cell] (joint : MeasureTheory.Measure (X × PKG25NoFreeLunch.Label))
      [MeasureTheory.IsProbabilityMeasure joint] (cell : X → Cell),
      Measurable cell →
        ∀ (A : Set ℝ),
          MeasurableSet A →
            ∫ (z : X × PKG25NoFreeLunch.Label),
                {z | PKG25NoFreeLunch.jointLawPartitionPredictor joint cell z.1 ∈ A ∧ z.2 = true}.indicator (fun x => 1)
                  z ∂joint =
              ∫ (z : X × PKG25NoFreeLunch.Label),
                {z | PKG25NoFreeLunch.jointLawPartitionPredictor joint cell z.1 ∈ A}.indicator
                  (fun z => PKG25NoFreeLunch.jointLawPartitionPredictor joint cell z.1) z ∂joint
\end{ReviewVerbatim}

\paragraph{Lean proof endpoint (built separately)}\begin{ReviewVerbatim}
PKG25NoFreeLunch.source_partition_finite_measurable_repairSpec_proof
\end{ReviewVerbatim}

\paragraph{Recorded source-to-Spec screening}
\textbf{Verdict:} \texttt{not recorded}
\\
\textbf{Reason:} not recorded
\reviewmatch{Matches source input}
\noindent\textbf{Reviewer annotation}\par
\reviewerbox
\clearpage
\hypertarget{source-claim-1b69269ab69dfcc0}{}
\section*{Review row 17}
\textbf{Source locator:} {\footnotesize\raggedright cited publication, lines 135-\allowbreak{}148\par}
\paragraph{Verbatim source input}
\begin{ReviewVerbatim}
\subsection{Main Proof}
In the remainder of this section, we prove \Cref{thm:main} in full. We first rewrite \Cref{thm:main} in an equivalent formulation.

\setcounter{theorem}{0}
\begin{theorem}
For a collaboration strategy $\CC$, $\acc_\CC(S) \ge \min_{i\in [n]}\acc_i(S)$ for all collaboration settings $S$ if and only if there exists $k\in [n]$ and $\alpha \in \{0,1\}$ such that for all $(p_1,p_2,\cdots,p_n)\in (0,1)^n$:

\begin{enumerate}
    \item[(i)] If $p_k\neq \frac{1}{2},$ then
    $\CC(p_1,p_2,\cdots,p_n)=\round{p_k}.$
    \item[(ii)] If $p_k=\frac{1}{2}$, then
        $\CC(p_1,p_2,\cdots,p_n) = \alpha.$
\end{enumerate}
\end{theorem}
\end{ReviewVerbatim}


\paragraph{Expanded PaperInterface specification}
\begin{ReviewVerbatim}
PKG25NoFreeLunch.NonCollaborative PKG25NoFreeLunch.boundaryFlipStrategy ∧
  ¬PKG25NoFreeLunch.ReliableJointLaw PKG25NoFreeLunch.boundaryFlipStrategy
\end{ReviewVerbatim}

\paragraph{Lean proof endpoint (built separately)}\begin{ReviewVerbatim}
PKG25NoFreeLunch.source_unnumbered_iff_restatementSpec_proof
\end{ReviewVerbatim}

\paragraph{Recorded source-to-Spec screening}
\textbf{Verdict:} \texttt{not recorded}
\\
\textbf{Reason:} not recorded
\reviewmatch{Matches source input}
\noindent\textbf{Reviewer annotation}\par
\reviewerbox
\clearpage
\hypertarget{source-claim-f706fa3b58e00679}{}
\section*{Review row 18}
\textbf{Source locator:} {\footnotesize\raggedright cited publication, lines 156-\allowbreak{}160\par}
\paragraph{Verbatim source input}
\begin{ReviewVerbatim}
To show \Cref{prop:part1}, suppose for sake of contradiction that there does not exist $k\in [n]$ satisfying this property. This means that for every $k\in [n]$, there must exist some tuple $(p_1,p_2,\cdots,p_n)\in (0,1)^n$ where $p_k\neq \frac{1}{2},$ such that $\CC(p_1,p_2,\cdots,p_n)\neq \round{p_k}.$ We show in  \Cref{lem:counterexample} below that the existence of such a tuple implies that there is a collaboration setting $S_k$ such that $\acc_k(S_k)>\acc_\CC(S_k)$ and $\acc_i(S_k)\ge \acc_\CC(S_k)$ for all $i\in [n]\setminus \{k\}.$ Since we can construct such an $S_k$ for all $k\in [n]$, \Cref{prop:linear-combination} implies the existence of a collaboration setting $S$ such that
\begin{equation}
     \acc_\CC(S) = \sum_{k=1}^n \frac{1}{n} \acc_\CC(S_k) < \sum_{k=1}^n \frac{1}{n} \acc_i(S_k) = \acc_i(S)
\end{equation}
for all $i\in [n],$ providing the desired contradiction. Therefore, it suffices to show \Cref{lem:counterexample}.
\end{ReviewVerbatim}


\paragraph{Expanded PaperInterface specification}
\begin{ReviewVerbatim}
∀ {n : ℕ} [Nonempty (Fin n)] (C : PKG25NoFreeLunch.CollaborationStrategy n),
  (¬∃ k, PKG25NoFreeLunch.DefersAwayFromHalf C k) →
    ∃ S,
      PKG25NoFreeLunch.source_assumption_strategy_expectation_well_formed S C ∧
        ∀ (i : Fin n), S.strategyAccuracy C < S.agentAccuracy i
\end{ReviewVerbatim}

\paragraph{Lean proof endpoint (built separately)}\begin{ReviewVerbatim}
PKG25NoFreeLunch.source_formula_proposition7_average_mixtureSpec_proof
\end{ReviewVerbatim}

\paragraph{Recorded source-to-Spec screening}
\textbf{Verdict:} \texttt{not recorded}
\\
\textbf{Reason:} not recorded
\reviewmatch{Matches source input}
\noindent\textbf{Reviewer annotation}\par
\reviewerbox
\clearpage
\hypertarget{source-claim-2c1cc7b2a1433aaa}{}
\section*{Review row 19}
\textbf{Source locator:} {\footnotesize\raggedright cited publication, lines 255-\allowbreak{}256\par}
\paragraph{Verbatim source input}
\begin{ReviewVerbatim}
Now set $\AA_k = \{\{0,1\}\}$ and $\AA_i = \{\{0\}, \{1\}\}$ for all $i\neq k.$ Then $P_k(0)=P_k(1)=\frac{2}{3} - \frac{\epsilon}{3}$, while $P_i(0)=\epsilon$ and $P_i(1)=1-\epsilon$ for $i\neq k.$ Then $\Yhat_\CC(0) = \Yhat_k(0)$ and $\Yhat_\CC(1) = \Yhat_k(1)$ since $(P_1(0), \cdots, P_n(0)), (P_1(1), \cdots, P_n(1))\in (0,1)^n.$
Then observe that $\acc_k(S_1)=\acc_\CC(S_1)=\frac{2}{3}-\frac{\epsilon}{3}$ (since $\CC$ always defers to agent $k$'s classification) and $\acc_i(S_1)=1 - \epsilon$ for $i\neq k$. For $\epsilon < \frac{1}{2}, \frac{2}{3}-\frac{\epsilon}{3} < 1 - \epsilon,$ giving the desired conclusion.
\end{ReviewVerbatim}


\paragraph{Expanded PaperInterface specification}
\begin{ReviewVerbatim}
∀ {n : ℕ} {ε : ℝ} (hε0 : 0 < ε) (hεhalf : ε < 1 / 2) {C : PKG25NoFreeLunch.CollaborationStrategy n} {k : Fin n},
  PKG25NoFreeLunch.DefersAwayFromHalf C k →
    PKG25NoFreeLunch.FiniteCollaborationSetting.strategyAccuracy C
          (PKG25NoFreeLunch.part2S1ParamSetting ε hε0 hεhalf k) =
        2 / 3 - ε / 3 ∧
      (PKG25NoFreeLunch.part2S1ParamSetting ε hε0 hεhalf k).agentAccuracy k = 2 / 3 - ε / 3 ∧
        ∀ (i : Fin n),
          i ≠ k →
            (PKG25NoFreeLunch.part2S1ParamSetting ε hε0 hεhalf k).agentAccuracy i = 1 - ε ∧
              PKG25NoFreeLunch.FiniteCollaborationSetting.strategyAccuracy C
                  (PKG25NoFreeLunch.part2S1ParamSetting ε hε0 hεhalf k) <
                (PKG25NoFreeLunch.part2S1ParamSetting ε hε0 hεhalf k).agentAccuracy i
\end{ReviewVerbatim}

\paragraph{Lean proof endpoint (built separately)}\begin{ReviewVerbatim}
PKG25NoFreeLunch.source_formula_proposition9_s1_accuraciesSpec_proof
\end{ReviewVerbatim}

\paragraph{Recorded source-to-Spec screening}
\textbf{Verdict:} \texttt{not recorded}
\\
\textbf{Reason:} not recorded
\reviewmatch{Matches source input}
\noindent\textbf{Reviewer annotation}\par
\reviewerbox
\clearpage
\hypertarget{source-claim-6680595ed06f984e}{}
\section*{Review row 20}
\textbf{Source locator:} {\footnotesize\raggedright cited publication, lines 302-\allowbreak{}311\par}
\paragraph{Verbatim source input}
\begin{ReviewVerbatim}
The result follows by applying \Cref{prop:linear-combination} with $S_1$ and $S_2$, taking $\lambda_1$ sufficiently close to $1$. In particular, \Cref{prop:linear-combination} implies that for all $\lambda$, there exists a collaboration setting $S$ such that
\begin{equation}
    \acc_i(S) = \lambda \acc_i(S_1) + (1-\lambda) \acc_i(S_2)
\end{equation}
for all $i\in [n]$, and
\begin{equation}
    \acc_\CC(S) = \lambda \acc_\CC(S_1) + (1-\lambda) \acc_\CC(S_2).
\end{equation}
Then, for all $\lambda < 1$, since $\acc_k(S_1)=\acc_\CC(S_1)$ and $\acc_k(S_2)>\acc_\CC(S_2),$ we have that $\acc_k(S) > \acc_\CC(S).$
Now, since $\acc_i(S_1)>\acc_\CC(S_1)$, regardless of $\acc_i(S_2), \acc_\CC(S_2),$ by taking $\lambda$ sufficiently close to $1$, $\acc_i(S) > \acc_\CC(S).$ This provides the desired contradiction.
\end{ReviewVerbatim}


\paragraph{Expanded PaperInterface specification}
\begin{ReviewVerbatim}
(∀ {n : ℕ} (S1 S2 : PKG25NoFreeLunch.FiniteCollaborationSetting n),
    ∃ Smix,
      (∀ (i : Fin n), Smix.agentAccuracy i = 7 / 8 * S1.agentAccuracy i + 1 / 8 * S2.agentAccuracy i) ∧
        ∀ (C : PKG25NoFreeLunch.CollaborationStrategy n),
          PKG25NoFreeLunch.source_assumption_strategy_expectation_well_formed Smix C ∧
            Smix.strategyAccuracy C =
              7 / 8 * PKG25NoFreeLunch.FiniteCollaborationSetting.strategyAccuracy C S1 +
                1 / 8 * PKG25NoFreeLunch.FiniteCollaborationSetting.strategyAccuracy C S2) ∧
  ∀ {n : ℕ} [Nonempty (Fin n)] {C : PKG25NoFreeLunch.CollaborationStrategy n} {k : Fin n} {p q : Fin n → ℝ},
    PKG25NoFreeLunch.DefersAwayFromHalf C k →
      PKG25NoFreeLunch.Interior p →
        PKG25NoFreeLunch.Interior q →
          p k = 1 / 2 →
            q k = 1 / 2 →
              C p = true →
                C q = false →
                  ∃ Smix,
                    PKG25NoFreeLunch.source_assumption_strategy_expectation_well_formed Smix C ∧
                      ∀ (i : Fin n), Smix.strategyAccuracy C < Smix.agentAccuracy i
\end{ReviewVerbatim}

\paragraph{Lean proof endpoint (built separately)}\begin{ReviewVerbatim}
PKG25NoFreeLunch.source_formula_proposition9_final_mixtureSpec_proof
\end{ReviewVerbatim}

\paragraph{Recorded source-to-Spec screening}
\textbf{Verdict:} \texttt{not recorded}
\\
\textbf{Reason:} not recorded
\reviewmatch{Matches source input}
\noindent\textbf{Reviewer annotation}\par
\reviewerbox

\end{document}
